\documentclass[11pt,letterpaper]{article}
\usepackage{fontspec}
\usepackage{tocloft}
\usepackage{multicol}
\usepackage{fancyhdr}
\usepackage{nicefrac}
\usepackage[left=1.3in,right=1.3in,top=1in,bottom=1in]{geometry}
\setmainfont[Scale=1.3,AutoFakeBold=1.5,AutoFakeSlant=0.3]{Doves Type}

\widowpenalty=10000
\clubpenalty=10000
\interlinepenalty=500

\title{Pumpkin-Pot Chicken}
\author{}
\date{}

\pagestyle{fancy}
\fancyhf{}
\fancyhead[C]{\textit{Pumpkin-Pot Chicken}}
\fancyfoot[C]{\thepage}
\renewcommand{\headrulewidth}{0pt}
\renewcommand{\footrulewidth}{0pt}

\begin{document}

\maketitle
\thispagestyle{empty}

 
%\begin{quote}
%\textit{Chicken leg quarters brown, then braise inside a covered cooking pumpkin with onion, leek, shallot, garlic, thyme, and sage. White wine and stock become the broth, and corn kernels go in near the end. Serve the soft pumpkin, the chicken, the corn, and the broth. Serves 4.}
%\end{quote}

\section*{Ingredients}
\setlength{\columnsep}{20pt}
\begin{multicols}{2}
\noindent
    Cooking pumpkin \dotfill 1 (8--12 lb.) \\
    Chicken leg quarters, skin-on \dotfill 4 (3--3\nicefrac{1}{2} lb.) \\
    Kosher salt \dotfill 1 Tbsp. plus more \\
    Black pepper \dotfill 1 tsp. \\
    Vegetable oil \dotfill 1 Tbsp. \\
    Yellow onion, large \dotfill 1 \\
    Leek, large \dotfill 1 \\
    Shallots \dotfill 2 \\
    \columnbreak
    Garlic cloves \dotfill 6 \\
    Dried thyme \dotfill 1 tsp. \\
    Rubbed sage \dotfill \nicefrac{3}{4} tsp. \\
    Bay leaves \dotfill 2 \\
    Dry white wine \dotfill \nicefrac{1}{2} cup \\
    Chicken stock \dotfill 1 cup \\
    Ears of corn \dotfill 2 \\
    Apple cider vinegar \dotfill 2 Tbsp. \\
    Unsalted butter \dotfill 1 Tbsp. \\
\end{multicols}

\section*{Directions}

\noindent
Heat the oven to \textit{325°F} ---
Wash the \textbf{pumpkin}. Cut a lid large enough for the \textbf{chicken}, with the knife angled so the lid sits down into the opening. Scrape out the seeds and strings. Salt the inside lightly ---
Dice the \textbf{onion}. Halve the \textbf{leek} lengthwise, rinse out the grit, and slice the white and pale green. Slice the \textbf{shallots}. Mince the \textbf{garlic} ---
Cut the kernels from the \textbf{corn}. Break the cobs in half ---
Pat the \textbf{chicken} dry. Season it all over with 1~Tbsp. \textbf{kosher~salt} and \nicefrac{1}{2}~tsp. \textbf{black~pepper}

\begin{enumerate}
    \item Set a large cast-iron skillet or other all-metal pan over medium-high heat. Add 1~Tbsp. \textbf{vegetable~oil}. Brown the \textbf{chicken} in two batches, skin-side down, without moving it, for \textit{8--10~minutes}, until the skin is deep golden and releases from the pan. Turn and cook \textit{2--3~minutes} more, until the other side has color. Move the \textbf{chicken} to a plate. Leave about 2~Tbsp. of fat in the pan.

    \item Lower the heat to medium. Add the \textbf{onion}, \textbf{leek}, and \textbf{shallots} with a pinch of \textbf{kosher~salt}. Cook, stirring now and then, until they soften and the pan goes from wet to pale gold, \textit{10--12~minutes}. Stir in the \textbf{garlic}, 1~tsp. \textbf{dried~thyme}, and \nicefrac{3}{4}~tsp. \textbf{rubbed~sage}. Cook for \textit{1~minute}, until the garlic smells sweet and stays pale.

    \item Add \nicefrac{1}{2}~cup \textbf{dry~white~wine}. Scrape the fond off the pan. Cook until only a glaze is left on the vegetables, \textit{2--4~minutes}. Add 1~cup \textbf{chicken~stock} and the \textbf{bay~leaves}. Bring it to a simmer.

    \item Set the \textbf{chicken} skin-side up in the \textbf{pumpkin}. Tuck the \textbf{corn} cobs around it where they fit. Pour the pan contents around the \textbf{chicken}. Set the lid on. Cover the stem with foil. Put the \textbf{pumpkin} in the pan and bake at \textit{325°F} for \textit{1\nicefrac{1}{2}~hours}.

    \item Lift the lid. Take out the cobs and the \textbf{bay~leaves}. Add the \textbf{corn} kernels. Cover again and bake \textit{30~minutes} more. It is done when the thickest part of a thigh reads \textit{185°F}, the meat has pulled back at the joint, and a knife slides into the \textbf{pumpkin} wall with little resistance. Continue in \textit{15~minute} stretches until it does. The skin will be soft.

    \item Rest, covered, for \textit{15~minutes}. Move the \textbf{chicken} to a platter. Skim the fat from the broth. Stir in 1~Tbsp. \textbf{unsalted~butter}. Add 1~Tbsp. \textbf{apple~cider~vinegar}, then more, up to the remaining tablespoon, until the broth tastes savory and bright. Season with the remaining \textbf{black~pepper} and more \textbf{kosher~salt}. Serve each bowl with \textbf{pumpkin} flesh scraped off the skin, \textbf{chicken}, \textbf{corn}, and broth. It is best in the first \textit{20~minutes}.
\end{enumerate}

\newpage

{\small
\setlength{\columnsep}{20pt}
\setlength{\multicolsep}{6pt}
\begin{multicols}{2}
\setlength{\parindent}{0pt}
\setlength{\parskip}{4pt}

\section*{Equipment Required}
\begin{itemize}
    \item Cast-iron skillet or other all-metal pan wide enough to hold the pumpkin and fit in the oven
    \item Large plate
    \item Tongs
    \item Wooden spoon
    \item Knife and cutting board
    \item Measuring cups and spoons
    \item Aluminum foil
    \item Instant-read thermometer
    \item Ladle
\end{itemize}

\section*{Hints and Notes}

\subsection*{Yield}
\begin{itemize}
    \item Serves 4.
    \item One quarter is a serving. The pumpkin and the broth fill the bowls.
\end{itemize}

\subsection*{Mise en Place}
\begin{itemize}
    \item Heat the oven to \textit{325°F} before you brown the \textbf{chicken}.
    \item Cut and scrape the \textbf{pumpkin} first. The lid has to fit back on.
    \item Cut the \textbf{onion}, \textbf{leek}, \textbf{shallots}, and \textbf{garlic}. Keep them on the board.
    \item Cut the \textbf{corn} kernels free. Break the cobs in half.
    \item Pat the \textbf{chicken} dry and season it.
    \item Measure the \textbf{wine}, \textbf{stock}, \textbf{thyme}, \textbf{sage}, \textbf{vinegar}, and \textbf{butter} before the pan is hot.
\end{itemize}

\subsection*{Ingredient Tips}
\begin{itemize}
    \item Use a cooking pumpkin, such as Cinderella or Long Island Cheese. A carving pumpkin is watery. A sugar pie pumpkin is too small for four quarters.
    \item The pumpkin and the pan have to fit in the oven together.
    \item Eight skin-on \textbf{drumsticks} can stand in for the quarters. Use the same heat and the same tests for doneness.
    \item \textbf{Kosher~salt} here is Diamond Crystal. Morton is denser. Use 2~tsp. Morton on the \textbf{chicken} in place of 1~Tbsp. Diamond Crystal.
    \item Use rubbed \textbf{sage}. Ground sage is stronger. Use \nicefrac{1}{2}~tsp. ground sage in its place.
    \item \textbf{Dry~white~wine}: Sauvignon Blanc or Pinot Grigio. Skip anything sweet.
    \item A fresh pumpkin gives off a lot of liquid. The 1~cup of \textbf{stock} is enough.
\end{itemize}

\subsection*{Preparation Tips}
\begin{itemize}
    \item Angle the lid cut so the top sits in the hole and the steam stays inside.
    \item Foil on the stem keeps it from scorching.
    \item Brown in two batches. A crowded pan steams the skin and leaves a pale fond.
    \item The \textbf{pumpkin} wall cooks more slowly than the meat. Dark meat can stay in until the wall yields to a knife.
    \item Take the cobs out when the kernels go in. They have already flavored the broth.
    \item A thin broth can go back into the pan and simmer \textit{5--8~minutes} before you add the \textbf{vinegar} and \textbf{butter}.
\end{itemize}

\subsection*{Make Ahead \& Storage}
\begin{itemize}
    \item Cook it the day you serve it. The \textbf{pumpkin} is at its best the first night.
    \item Refrigerate leftover \textbf{chicken}, \textbf{pumpkin} flesh, and broth up to \textit{3~days}.
    \item Reheat, covered, over low heat until the \textbf{chicken} is hot.
    \item Blend leftover \textbf{pumpkin} flesh with leftover broth for soup. Warm it gently and taste for salt and \textbf{vinegar}.
\end{itemize}

\subsection*{Serving Suggestions}
\begin{itemize}
    \item Bring the \textbf{pumpkin} to the table in the pan. Ladle the broth, then lift out the \textbf{chicken}.
    \item Crusty bread is the right thing for the broth.
    \item It stays hot in the \textbf{pumpkin} after the first \textit{20~minutes}. The skin stays soft.
\end{itemize}

\end{multicols}
}

\end{document}
